\textbf{Phần 1. Trắc nghiệm lựa chọn}
\setcounterpageref{ex}{0}

\begin{ex} %[2D14V71]
Tính giá trị của biểu thức: $5 \times 8 + 15$.
\choice
{$45$}
{$50$}
{$60$}
{\True $55$}
\loigiai{
$5 \times 8 = 40$; $40 + 15 = 55$.
\textbf{Đáp án của câu hỏi là D}
}
\end{ex}

\begin{ex} %[2D14V71]
Tính giá trị của biểu thức: $45 : 5 + 28$.
\choice
{$35$}
{\True $37$}
{$39$}
{$47$}
\loigiai{
$45 : 5 = 9$; $9 + 28 = 37$.
\textbf{Đáp án của câu hỏi là B}
}
\end{ex}

\begin{ex} %[2D14V71]
Tìm $x$ biết: $x \times 2 = 18$.
\choice
{$8$}
{$16$}
{\True $9$}
{$36$}
\loigiai{
$x = 18 : 2 = 9$.
\textbf{Đáp án của câu hỏi là C}
}
\end{ex}

\begin{ex} %[2D14V71]
Có $35$ chiếc kẹo chia đều cho $5$ bạn. Mỗi bạn nhận được bao nhiêu chiếc kẹo?
\choice
{\True $7$ chiếc}
{$6$ chiếc}
{$8$ chiếc}
{$5$ chiếc}
\loigiai{
$35 : 5 = 7$ (chiếc kẹo).
\textbf{Đáp án của câu hỏi là A}
}
\end{ex}

\begin{ex} %[2D14V71]
Một túi gạo cân nặng $5\text{ kg}$. Hỏi $9$ túi gạo như thế cân nặng bao nhiêu ki-lô-gam?
\choice
{$40\text{ kg}$}
{$50\text{ kg}$}
{$14\text{ kg}$}
{\True $45\text{ kg}$}
\loigiai{
$5 \times 9 = 45$ (kg).
\textbf{Đáp án của câu hỏi là D}
}
\end{ex}

\begin{ex} %[2D14V71]
Tìm $x$ biết: $x : 5 = 6$.
\choice
{$25$}
{\True $30$}
{$35$}
{$11$}
\loigiai{
$x = 6 \times 5 = 30$.
\textbf{Đáp án của câu hỏi là B}
}
\end{ex}

\begin{ex} %[2D14V71]
Tính: $2 \times 7 + 5 \times 4$.
\choice
{$32$}
{$36$}
{\True $34$}
{$38$}
\loigiai{
$2 \times 7 = 14$; $5 \times 4 = 20$; $14 + 20 = 34$.
\textbf{Đáp án của câu hỏi là C}
}
\end{ex}

\begin{ex} %[2D14V71]
Một cuộn dây dài $20\text{ dm}$ được cắt thành các đoạn dài $2\text{ dm}$. Cắt được bao nhiêu đoạn dây?
\choice
{\True $10$ đoạn}
{$8$ đoạn}
{$12$ đoạn}
{$5$ đoạn}
\loigiai{
$20 : 2 = 10$ (đoạn).
\textbf{Đáp án của câu hỏi là A}
}
\end{ex}

\begin{ex} %[2D14V71]
Điền dấu thích hợp: $5 \times 6 \dots 2 \times 10 + 10$.
\choice
{$>$}
{$<$}
{Không so sánh được}
{\True $=$}
\loigiai{
$5 \times 6 = 30$; $2 \times 10 + 10 = 30$. Hai vế bằng nhau.
\textbf{Đáp án của câu hỏi là D}
}
\end{ex}

\begin{ex} %[2D14V71]
Có $16$ học sinh xếp thành các hàng, mỗi hàng có $2$ bạn. Xếp được bao nhiêu hàng?
\choice
{$7$ hàng}
{\True $8$ hàng}
{$9$ hàng}
{$6$ hàng}
\loigiai{
$16 : 2 = 8$ (hàng).
\textbf{Đáp án của câu hỏi là B}
}
\end{ex}

\begin{ex} %[2D14V71]
Một người nuôi $5$ chuồng thỏ, mỗi chuồng có $4$ con thỏ. Người đó bán đi $6$ con thỏ. Còn lại bao nhiêu con thỏ?
\choice
{$12$ con}
{$16$ con}
{\True $14$ con}
{$20$ con}
\loigiai{
Tổng số thỏ: $5 \times 4 = 20$ con. Còn lại: $20 - 6 = 14$ con.
\textbf{Đáp án của câu hỏi là C}
}
\end{ex}

\begin{ex} %[2D14V71]
Tính: $0 : 5 + 2 \times 9$.
\choice
{\True $18$}
{$0$}
{$19$}
{$23$}
\loigiai{
$0 : 5 = 0$; $2 \times 9 = 18$; $0 + 18 = 18$.
\textbf{Đáp án của câu hỏi là A}
}
\end{ex}

\textbf{Phần 2. Trắc nghiệm đúng sai}
\setcounterpageref{ex}{0}

\begin{ex} %[2D14V71]
Xét tính đúng sai của các phép tính nhân, chia:
\choiceTF
{\True Số 0 nhân với bất kì số nào cũng bằng 0.}
{Số 0 chia cho số nào cũng bằng chính số đó.}
{\True 5 x 7 = 35 và 35 : 5 = 7.}
{\True Nếu tích là 20 và một thừa số là 2 thì thừa số kia là 10.}
\loigiai{
\begin{itemize}
    \item Ý a) Đúng: Đúng tính chất số 0 trong phép nhân.
    \item Ý b) Sai: 0 chia cho số khác 0 luôn bằng 0, không phải chính số đó.
    \item Ý c) Đúng: Quan hệ giữa phép nhân và phép chia. Đúng.
    \item Ý d) Đúng: 20 : 2 = 10. Đúng.
\end{itemize}
\begin{itemize}
    \item a) \textbf{Đúng}.
    \item b) \textbf{Sai}.
    \item c) \textbf{Đúng}.
    \item d) \textbf{Đúng}.
\end{itemize}
}
\end{ex}

\begin{ex} %[2D14V71]
Một nhóm gồm 6 bạn, mỗi bạn góp 5 quyển vở. Số vở này đem chia đều cho 3 tổ học sinh nghèo:
\choiceTF
{\True Tổng số vở các bạn góp là 30 quyển.}
{\True Mỗi tổ học sinh nghèo nhận được 10 quyển vở.}
{Nếu mỗi bạn góp thêm 1 quyển thì tổng số vở là 32 quyển.}
{\True Số vở của 2 bạn bằng 1/3 tổng số vở cả nhóm.}
\loigiai{
\begin{itemize}
    \item Ý a) Đúng: 6 x 5 = 30 quyển. Đúng.
    \item Ý b) Đúng: 30 : 3 = 10 quyển. Đúng.
    \item Ý c) Sai: 6 x 6 = 36 quyển, không phải 32. Sai.
    \item Ý d) Đúng: 2 x 5 = 10 quyển = 30 : 3. Đúng.
\end{itemize}
\begin{itemize}
    \item a) \textbf{Đúng}.
    \item b) \textbf{Đúng}.
    \item c) \textbf{Sai}.
    \item d) \textbf{Đúng}.
\end{itemize}
}
\end{ex}

\textbf{Phần 3. Trả lời ngắn}
\setcounterpageref{ex}{0}

\begin{ex} %[2D14V71]
Tính giá trị: $5 \times 9 - 18$.
\shortans{27}
\loigiai{
\textbf{Đáp số:} $27$\\
$45 - 18 = 27$.
}
\end{ex}

\begin{ex} %[2D14V71]
Tìm $x$ biết: $x : 2 = 8$.
\shortans{16}
\loigiai{
\textbf{Đáp số:} $16$\\
$x = 8 \times 2 = 16$.
}
\end{ex}

\begin{ex} %[2D14V71]
Có 40 quả cam xếp đều vào 5 đĩa. Hỏi mỗi đĩa có bao nhiêu quả cam?
\shortans{8}
\loigiai{
\textbf{Đáp số:} $8$\\
$40 : 5 = 8$ quả.
}
\end{ex}

\begin{ex} %[2D14V71]
Mỗi can chứa 5 lít dầu. Hỏi 7 can như thế chứa bao nhiêu lít dầu?
\shortans{35}
\loigiai{
\textbf{Đáp số:} $35$\\
$5 \times 7 = 35$ lít.
}
\end{ex}

\begin{ex} %[2D14V71]
Một sợi dây dài 18 dm cắt thành các đoạn dài 2 dm. Hỏi cắt được bao nhiêu đoạn?
\shortans{9}
\loigiai{
\textbf{Đáp số:} $9$\\
$18 : 2 = 9$ đoạn.
}
\end{ex}

\begin{ex} %[2D14V71]
Tính: $2 \times 8 : 4$.
\shortans{4}
\loigiai{
\textbf{Đáp số:} $4$\\
$16 : 4 = 4$.
}
\end{ex}
